\documentclass[11pt,letterpaper]{article}
\usepackage[margin=1in]{geometry}
\usepackage[T1]{fontenc}
\usepackage{lmodern}
\usepackage{array,booktabs,tabularx}
\setlength{\parindent}{0pt}
\setlength{\parskip}{7pt}
\newcommand{\code}[1]{\texttt{#1}}
% Four planned pages: architecture; TCP; tests; AI use and collaboration.
% Technical results describe implementation commit 637aebb.
\begin{document}
\begin{center}
{\Large\bfseries COMP512 Assignment 1: Travel Reservation System}\\[6pt]
Group 30 --- Zuojun Gu and Yinkun Zhou\\
McGill University, Fall 2026
\end{center}

\section{Our design and RMI version}
We split the application into a console client, middleware, and three resource managers:
one each for flights, cars, and rooms. Both RMI and TCP use the same arrangement.
The client talks to middleware, which sends inventory requests to the right manager.

We keep customers at middleware. This makes it easier to build one bill containing
flights, cars, and rooms, without copying customer records between servers. Each manager
keeps its own inventory and booking records. Customers, bundles, and deletion are handled
by shared middleware logic, so the two transports follow the same rules.

\subsection{RMI}
The RMI client looks up middleware in the registry and calls the original
\code{IResourceManager} interface. Middleware looks up the three manager stubs and
forwards requests. We kept the console commands and interface, with small starter fixes
for compilation, bill totals, and bundle flags. Hosts, ports, and registry prefixes
can be configured for deployment.

We added two internal inventory methods: \code{reserveInventory} and
\code{releaseInventory}. Reserving returns a receipt with the quantity and price, and
uses a booking ID so repeating a release does not restore inventory twice.
Middleware uses these methods without creating customers at the managers.

\subsection{Customers and bundles}
For a bundle, we first check the flight inputs and count repeated flight numbers.
We then reserve the flights and any requested car or room. The customer's bill is
updated only after everything succeeds. If a later item is unavailable, we cancel
this attempt's bookings and leave earlier reservations alone.

Deleting a customer releases their bookings at all three managers. If cleanup fails,
we keep the unfinished work so it can be retried while the endpoints remain usable.
Billing follows the starter's grouped latest-price rule: a later booking at a new price
can change the unit price shown for earlier bookings of the same item.

\subsection{Which RMI version we chose}
Zuojun developed the AI-assisted version, while Yinkun independently developed the
other version without AI. We compared them and chose the AI version because its error
and exception handling was more detailed, which we thought would help avoid potential
bugs. For example, it keeps track of unfinished cleanup and interrupted deletion.
We saved the accepted version under the Git tag \code{rmi-complete}; the other
preliminary version is a separate submission attachment.

\newpage
\section{TCP version}
\subsection{One message service for all calls}
We use TCP sockets for both client--middleware and middleware--manager communication.
The console remains blocking: it sends a command and waits for its reply.

The client does not build a different message format for every command.
\code{TcpChannel.invoke(Method, Object[])} handles general calls, and
\code{Services.blocking} wraps it as the normal client interface.
All client methods, including both \code{newCustomer} versions, go through this service.
\code{Protocol.read/write} handles encoding on both links.

\begin{center}
\small
\begin{tabularx}{\linewidth}{@{}lX@{}}
\toprule
Message & Main fields\\
\midrule
Request & Length, version, kind, request ID, method signature, typed arguments\\
Response & Length, version, kind, request ID, error or typed result\\
\bottomrule
\end{tabularx}
\end{center}
Length-prefixed UTF-8 strings support multiline bills and location names. Frame lengths
handle split or combined TCP reads. We reject invalid methods, types, and lengths,
with a 1 MiB limit. Both links use our binary codec, without RMI or Java serialization.

\subsection{Keeping middleware responsive}
Each connection has a reader and writer. The channel assigns each request an ID and
keeps a future for its reply. When a response arrives, its ID tells us which future
to complete, even if replies arrive out of order.

Middleware forwards requests and uses future callbacks to queue replies, without
waiting on backend futures. Socket threads can block on I/O, but unrelated requests
continue while a backend is slow.

Requests for the same customer form an ordered future chain. This avoids losing
reservations when two calls update one customer's bill. Other customers and inventory
queries can still progress. Resource managers accept multiple connections and use
four workers with a bounded queue. Inventory updates are synchronized so two requests
cannot both take the last seat. Queues and pending requests are bounded to limit overload.

\subsection{Errors and limitations}
A timeout (30 seconds by default) or disconnect closes the channel and fails waiting
calls. There is no automatic reconnect or resend: a mutation may already have happened.
Booking IDs prevent duplicate cleanup, and failed cancellations remain recorded for retry.

State is in memory and lost on restart; middleware is a single point of failure.
Bundle cleanup is not crash-safe, and queries may see temporary reservations.
Persistence and replication are outside our implementation.

\newpage
\section{What we tested}
All 14 suites passed on JDK 17.0.11 on October 3, 2026. We checked return values,
bills, and available quantities through public interfaces, real RMI endpoints, and
TCP sockets. These are the main cases we tried:

{\small
\renewcommand{\arraystretch}{1.2}
\begin{tabularx}{\linewidth}{@{}p{0.23\linewidth}X@{}}
\toprule
Case & Steps and result\\
\midrule
Add and reserve & Add two seats at \$100, then one at price 0. Query returns three seats at \$100. Reserve one: the bill is \$100, and deleting the reserved flight fails.\\
Last-seat race & Twelve callers try to reserve one seat. Exactly one succeeds. Releasing the same booking twice restores only one seat. Cancellation before a delayed reserve prevents that reserve.\\
Customers & Reject duplicate customers; combine all three resource types in one bill. Delete the customer and check that all quantities return. Backends have no corresponding customer record.\\
Successful bundle & Start with three \$100 seats, a \$30 car, and an \$80 room. Book two copies of the flight plus car and room: the bill is \$310 and one seat remains.\\
Failed bundle & Make a later item sold out. Earlier acquisitions from this attempt are released, while existing customer bookings stay intact. Empty or invalid bundles fail.\\
Errors and cleanup & Inject an error after a reserve changes inventory, then make cancellation fail. Retry after the fault clears and check for no leaked or double-restored inventory. Also resume interrupted deletion.\\
TCP messages & Round-trip all signatures, both customer overloads, Unicode, multiline bills, and receipts. Split and combine frames; reject oversized frames and unknown/internal client methods.\\
Concurrent requests & Hold 40 TCP flight requests. Cars queries on the same and another connection finish before Flights is released. Twelve bookings for one customer give a \$120 bill and are all released. RMI also checks independent progress and 100 distinct generated IDs.\\
Timeout/disconnect & Leave a TCP reply unfinished or close the server. Waiting calls fail instead of hanging indefinitely.\\
Separate processes & Run the console, middleware, and three managers in five JVMs for both transports. Test commands, prices, reservations, bills, bundles, deletion, and restored inventory. Exercise the Bash launchers too.\\
\bottomrule
\end{tabularx}}

Run the suite with \code{./tests/run-tests.ps1} or \code{bash tests/run-tests.sh}.
The Make build and Bash launcher checks also passed. Suite names and commands are
listed in \code{docs/TESTING.md}.

Yinkun ran the real local process tests and verified the final code. All local process
tests passed. We have not yet tested on five different machines, so the five-JVM result
only confirms the local process setup, not the required physical-machine demo.

\newpage
\section{AI use}
Zuojun used Codex for implementation, automated tests, documentation, and help drafting
this report. CodeGraph helped navigate the code. We kept stage commits and tested the
generated code rather than assuming it worked. Only this Codex conversation was used;
the complete logs are attached as \code{AI-conversations.json}.

The exported implementation logs record about \textbf{14.76 million tokens}
(14,763,681): 14,678,090 input and 85,591 output. Most input was cached
(14,190,976 tokens), so this is processed context rather than unique text we wrote.
We summed per-response usage once, without adding cumulative counters again.
The exports contain the main thread and two automatic-review threads from the same
conversation. The estimate covers implementation through October 3, 08:12:58 UTC
and excludes later report drafting. Details are in \code{AI-token-summary.json}.

\section{Our contributions and collaboration}
\begin{center}
\small
\renewcommand{\arraystretch}{1.2}
\begin{tabularx}{\linewidth}{@{}p{0.18\linewidth}XX@{}}
\toprule
Task & Zuojun Gu & Yinkun Zhou\\
\midrule
RMI & Independently built the AI version and discussed the final choice. & Independently built the non-AI version, compared versions, and tested and checked the final code.\\
TCP & Built the protocol, general wrapping, asynchronous middleware, and automated tests with AI. & Ran real local tests and checked the final code.\\
Customers / bundles & Implemented customer management, bills, bundle cleanup, and deletion with AI. & Tested and checked these features.\\
Report & Co-wrote the report, with AI drafting assistance. & Co-wrote the report and contributed code/test verification.\\
\bottomrule
\end{tabularx}
\end{center}

We first met on October 2 at 21:45 for about half an hour to read the assignment and
agree on the division of work. We met again on October 3 at around 03:00 for about
half an hour to compare the RMI versions and choose the AI-assisted one. Both times
are Montreal local time. For the later work, Zuojun focused on AI-assisted development,
while Yinkun focused on running tests and checking the final code. We wrote the report
together.
\end{document}
